\section{Enhetstester}
\label{sec:enhetstester}

Enhetstestene var skrevet med xUnit og brukte Moq for mocking av
avhengigheter. HTTP-laget ble mocket med en lokal
\texttt{StubHttpMessageHandler}-implementasjon som tok en lambda av typen
\texttt{Func<HttpRequestMessage, Task<HttpResponseMessage>>} og lot
testen kontrollere både request-inspeksjon og response-injeksjon uten
nettverkskall.

Komponentene som hadde dedikerte enhetstester var valgt ut fra to
kriterier: de inneholdt forretningslogikk som lot seg isolere fra
ASP.NET Cores request-pipeline, og de hadde feilstier som var krevende
å treffe gjennom integrasjonstester. Provider-laget
(\texttt{AntMediaProvider}, \texttt{NHNVideoProvider}) og
autentiseringshåndtererne (\texttt{AntMediaAuthHandler},
\texttt{NHNAuthHandler}) falt klart innenfor begge kriteriene.
\texttt{VideoProviderFactory}, helsesjekker, \texttt{GlobalExceptionHandler},
DI-registreringer og de egendefinerte exception-typene ble enhetstestet
av samme grunn. \texttt{VideoMetrics} ble bevisst utelatt fra
enhetstestene fordi klassen instrumenterer Prometheus via statiske
metoder uten testbar logikk; den ble i stedet verifisert indirekte
gjennom \gls{e2e}-fase~9 og manuell Grafana-sjekk
(jf.\ delkapittel~\ref{sec:e2e-testing}).

Et eksempel på en happy path-test verifiserte at
\texttt{CreateRoomAsync} sendte riktig request-body og mappet responsen
korrekt:

\begin{lstlisting}[language={[Sharp]C},
  caption={CreateRoomAsync — happy path}]
var result = await sut.CreateRoomAsync(
    new CreateRoomRequest { Provider = "antmedia", Name = "Demo Room" });

Assert.Equal("room-1", result.Id);
Assert.Contains("\"type\":\"liveStream\"", requestBody);
\end{lstlisting}

Feilscenarier ble testet like systematisk. Følgende test verifiserte at
en 500-respons fra AntMedia kastet \texttt{ExternalServiceException}
med original statuskode bevart, slik at \texttt{GlobalExceptionHandler}
kunne eksponere den i \texttt{externalStatusCode}-feltet i
ProblemDetails-responsen og la klienten skille AntMedia-feil fra
NHN-feil bak en felles \texttt{502 Bad Gateway}:

\begin{lstlisting}[language={[Sharp]C},
  caption={StopStreamAsync — upstream 500 kaster ExternalServiceException}]
var ex = await Assert.ThrowsAsync<ExternalServiceException>(
    () => sut.StopStreamAsync("room-1"));

Assert.Equal(500, ex.StatusCode);
\end{lstlisting}

\texttt{NHNVideoProviderTests.UnsupportedOperations\_ThrowNotSupportedException}
verifiserte at start, stopp og romtoken for NHN kastet
\texttt{NotSupportedException}.
\texttt{VideoProviderFactoryTests.GetProvider\_MixedCaseKey\_ReturnsRegisteredProvider} bekreftet case-insensitivt provider-oppslag med \texttt{''NhN''} som testdata.

\subsection{Token-håndtering}
\label{subsec:token-haandtering}

Autentiseringshåndtererne håndterte to sikkerhetskritiske operasjoner:
caching av token mellom kall og fornyelse når det cachete token ikke
lenger var gyldig. Én test verifiserte at et utløpt token utløste ny
henting ved at \texttt{NHNTokenCache.Expiry} ble satt til en tid i
fortiden. En separat test verifiserte at påfølgende kall med gyldig
cachet token kun utløste én token-henting.
\texttt{AntMediaAuthHandler} hadde tilsvarende cache- og signaturtester
for det lokalt signerte HMAC-SHA256-tokenet. Disse testene var en av
grunnene til at \texttt{NHNAuthHandler} (~97\,\%) og
\texttt{AntMediaAuthHandler} (~97\,\%) hadde høy linjedekning.